% !TEX encoding = UTF-8 Unicode
% BigDataSmallPrice -- Abschlussbericht
% ZHAW Big Data Seminar HS 2026
% Vorlage: ACM sigconf (zweispaltig)
\documentclass[sigconf]{acmart}

% ---- Sprache & Encoding ----
\usepackage[ngerman]{babel}
\usepackage[utf8]{inputenc}
\usepackage[T1]{fontenc}

% ---- Mathematik ----
\usepackage{amsmath}

% ---- Tabellen & Abbildungen ----
\usepackage{booktabs}
\usepackage{graphicx}
\usepackage{multirow}

% ---- TikZ ----
\usepackage{tikz}
\usetikzlibrary{shapes.geometric,arrows.meta,positioning,backgrounds,fit}

% ---- Plots ----
\usepackage{pgfplots}
\pgfplotsset{compat=1.18}
\usepgfplotslibrary{fillbetween}

% ---- Diverses ----
\usepackage{microtype}
\usepackage{url}
\usepackage{balance}
\usepackage{xcolor}
\emergencystretch=2em
\hyphenation{Zeit-rei-hen-da-ten-bank Zeit-rei-hen-da-ten-ban-ken
  Trainings-ab-fragen Ver-ar-bei-tungs-werk-zeuge Netz-nutzungs-tarife
  Wahr-schein-lich-keits-last-prog-no-sen Re-gu-la-ri-sie-rungs-pa-ra-me-ter
  Hoch-ver-fueg-bar-keits-in-fra-struk-tur}

% ---- ACM-Metadaten (Seminararbeit) ----
\setcopyright{none}
\copyrightyear{2026}
\acmYear{2026}
\acmDOI{}
\acmISBN{}
\acmConference[Big Data]{ZHAW Big Data}{HS26}{Winterthur, Schweiz}
\settopmatter{printacmref=false, printccs=false, printfolios=true}
\renewcommand\footnotetextcopyrightpermission[1]{}

% =========================================================
\begin{document}
% =========================================================

\title{BigDataSmallPrice: Automatisierte Prognose dynamischer Stromtarife
       f\"ur Winterthur mittels Machine Learning}

\author{Ryan Bachmann}
\affiliation{%
  \institution{ZHAW School of Engineering}
  \city{Winterthur}
  \country{Schweiz}}

\author{Dinis Silva Miguel}
\affiliation{%
  \institution{ZHAW School of Engineering}
  \city{Winterthur}
  \country{Schweiz}}

\author{Gian Ruchti}
\affiliation{%
  \institution{ZHAW School of Engineering}
  \city{Winterthur}
  \country{Schweiz}}

\renewcommand{\shortauthors}{Ryan Bachmann, Dinis Silva Miguel, Gian Ruchti}

% ---- Abstract ----
\begin{abstract}
Mit den neuen Schweizer Strommarktregeln werden dynamische Tarifmodelle ab 2026
f\"ur Netzbetreiber und Endkunden deutlich relevanter.
F\"ur Haushalte und Unternehmen bedeutet das eine grosse Chance: Wer weiss, wann
Strom g\"unstig ist, kann Waschmaschine, W\"armepumpe oder E-Auto gezielt zu
diesen Zeiten betreiben und bis zu 40\,\% Kosten sparen.
Das Problem: Bisher gibt es keine \"offentlich zug\"angliche Prognose, die mehrere
Tage im Voraus zeigt, wann der Strom billig oder teuer wird.

Wir l\"osen dieses Problem mit \emph{BigDataSmallPrice}, einer vollst\"andig
automatisierten Plattform, die t\"aglich Daten aus acht Quellen (darunter
europ\"aische Energieb\"orsendaten, Wetterprognosedaten und den historischen
Lastgang der Stadt Winterthur) zusammenf\"uhrt und daraus eine
7-Tage-Tarifprognose im 15-Minuten-Raster berechnet.
Zwei XGBoost-Modelle prognostizieren jeweils den Netzpreis (basierend auf der
lokalen Stromnachfrage) und den Energiepreis (basierend auf den europ\"aischen
B\"orsenpreisen).
Mit einem mittleren Fehler von~3{,}99\,\% beim Lastmodell und~7{,}72\,\%
beim Preismodell liegen beide deutlich unter den gesetzten Qualit\"atszielen.
Die Prognosen werden \"uber ein interaktives Dashboard mit Ampelindikator und
pers\"onlichen Verbrauchsempfehlungen angezeigt.
\end{abstract}

\maketitle
\pagestyle{fancy}
\fancyhf{}
\fancyhead[L]{Big Data, HS26, Winterthur, Schweiz}
\fancyhead[R]{Ryan Bachmann, Dinis Silva Miguel, Gian Ruchti}
\fancyfoot[C]{\thepage}


% =========================================================
\section{Einleitung}
% =========================================================

Die Schweizer Energieversorgung befindet sich im Umbruch.
Mit dem revidierten Stromversorgungsgesetz und der ab 2026 wirksamen
Tarifregulierung werden dynamische Netznutzungstarife regulatorisch
erm\"oglicht und von ersten Versorgern als Wahltarif eingef\"uhrt~\cite{ElCom2023,EKZ2024}.
Gleichzeitig steigt die Zahl steuerbarer Ger\"ate in Schweizer Haushalten rasant:
Der Marktanteil von Elektrofahrzeugen \"uberstieg 2024 die 25-Prozent-Marke,
und W\"armepumpen ersetzen zunehmend Heiz\-\"ol\-anlagen~\cite{ASTRA2024}.
Das Zusammenspiel von dynamischen Tarifen und steuerbaren Lasten er\"offnet
ein erhebliches Einsparpotenzial von bis zu 40\,\% gegen\"uber dem Fixpreis,
wenn der Verbrauch in Schwachlastzeiten verschoben wird~\cite{EKZ2024}.

Die Herausforderung besteht darin, dass Endkunden diese Einsparungen nur dann
gezielt nutzen k\"onnen, wenn sie \emph{im Voraus} wissen, wann Strom g\"unstig
oder teuer sein wird.
Bisherige Ans\"atze zur Strompreisvorhersage aus der Forschung~\cite{Lago2021,
Weron2014} fokussieren meist auf europ\"aische B\"orsenpreise, nicht aber auf
die konkrete Tarifzusammensetzung, die Endkunden auf der Rechnung sehen.
Kommerzielle Portale wie jenes der EKZ~\cite{EKZ2024} zeigen historische und
aktuelle Tarife, liefern aber keine Prognosen f\"ur die kommenden Tage.

\paragraph{Beitr\"age dieser Arbeit.}
Wir pr\"asentieren \emph{BigDataSmallPrice} mit folgenden Beitr\"agen:
\begin{enumerate}
  \item Eine vollautomatische Datenpipeline, die t\"aglich Daten aus acht Quellen
        einliest, bereinigt und in einer Zeitreihendatenbank speichert.
  \item Zwei kalibrierte Vorhersagemodelle f\"ur Netz- und Energiepreis mit
        nachvollziehbarem Retraining und messbaren Genauigkeitszielen.
  \item Transparente Tarifformeln, die Rohergebnisse in konkrete Rp./kWh-Werte
        mit Unsicherheitsgrenzen umrechnen.
  \item Ein Nutzerdashboard mit Wochenprognose, Ampelindikator und ger\"atespezifischen
        Empfehlungen sowie eine offene REST-Schnittstelle f\"ur Drittsysteme.
\end{enumerate}

% =========================================================
\section{Hintergrund und Problemstellung}
% =========================================================

\subsection{Dynamische Stromtarife in Winterthur}

Die EKZ (Elektrizit\"atswerke des Kantons Z\"urich) ver\"offentlichen seit 2024
zwei dynamische Tarifprodukte im 15-Minuten-Raster: einen \emph{Netzpreis}
(abh\"angig davon, wie stark das lokale Stromnetz gerade belastet ist) und einen
\emph{Energiepreis} (gekoppelt an den EPEX-Spot-Day-Ahead-Markt, auf dem
europ\"aische Stromproduzenten und -verbraucher handeln).
Der Gesamttarif, den Kunden zahlen, ergibt sich als Summe:
\begin{equation}
  T_{\mathrm{gesamt}}(t) = T_{\mathrm{netz}}(t) + T_{\mathrm{energie}}(t)
  \quad [\text{Rp./kWh}]
  \label{eq:tariff}
\end{equation}
In der Praxis schwankt dieser Gesamttarif zwischen rund 12 und 28~Rp./kWh,
ein Faktor von \"uber zwei zwischen gutem und schlechtem Zeitpunkt.
EKZ ver\"offentlicht die Tarife in der Regel erst am Vortag, was f\"ur eine
mehrt\"agige Verbrauchsplanung nicht ausreicht.

\subsection{Datengrundlage}

Tabelle~\ref{tab:datasources} fasst alle acht integrierten Datenquellen zusammen.
Die Grundlage f\"ur die Lastprognose bildet der Winterthurer Bruttolastgang
(OGD~\#1863): 15-min\"utige Strommessungen seit 2013, insgesamt rund 462\,000
Datenpunkte.
Die Energiepreisvorhersage basiert auf ENTSO-E-Day-Ahead-Preisen seit 2015.
Als Wetterdaten nutzen wir die kostenlose Open-Meteo-API f\"ur Winterthur und
zwei deutsche Referenzstationen, da norddeutsches Windaufkommen erfahrungs\-gem\"ass
den gr\"ossten Einfluss auf europ\"aische Strompreise hat~\cite{Lago2021}.
EKZ-, CKW- und Groupe-E-Tarife werden t\"aglich abgerufen und dienen als
Vergleichsreferenz zur Kalibrierung und Validierung.

\begin{table}[t]
  \caption{Integrierte Datenquellen}
  \label{tab:datasources}
  \footnotesize
  \setlength{\tabcolsep}{2.4pt}
  \begin{tabular}{@{}p{2.75cm}p{2.15cm}p{0.95cm}p{1.15cm}@{}}
    \toprule
    Quelle & Inhalt & Takt & Format \\
    \midrule
    ENTSO-E (A44/A65/\newline A75/A11/A01)
      & Preise, Last, Erzeug., Fl\"usse & 15\,min /\newline 1\,h & REST /\newline XML \\[2pt]
    Open-Meteo (CH + DE)
      & Temp., Wind, Solar, Wolken & 1\,h & REST /\newline JSON \\[2pt]
    OGD Winterthur \#1863
      & Bruttolastgang (ab 2013) & 15\,min & CSV /\newline API \\[2pt]
    OGD Winterthur \#3122
      & PV-Einspeisung & 15\,min & CSV /\newline API \\[2pt]
    EKZ-Tarif-API
      & Referenztarif (Netz+Energie) & 15\,min & REST /\newline JSON \\[2pt]
    CKW-Tarif-API
      & Vergleichstarif Innerschweiz & 15\,min & REST /\newline JSON \\[2pt]
    Groupe~E-API
      & Vergleichstarif Romandie & 15\,min & REST /\newline JSON \\[2pt]
    BAFU Hydrologie
      & Reservoir, Durchfluss & t\"aglich & CSV \\
    \bottomrule
  \end{tabular}
\end{table}

\subsection{Explorative Datenanalyse}

Abbildung~\ref{fig:eda} zeigt zwei zentrale Muster aus den Rohdaten.
\textbf{Links:} Das mittlere Lastprofil pro Stunde des Tages unterscheidet
sich deutlich zwischen Werktagen und Wochenenden.
An Werktagen entstehen zwei ausgep\"agte Spitzen: eine morgens (7--9~Uhr)
und eine abends (18--21~Uhr), w\"ahrend die Last zwischen 1 und 5~Uhr
auf rund 55~\% des Tagesmaximums f\"allt.
Am Wochenende ist das Profil flacher und die Abendspitze weniger stark.
Diese regelm\"assigen Muster sind der Hauptgrund, warum Zeitmerkmale
(Stunde, Wochentag, Kalender\-feier\-tage) die wichtigsten Eingabevariablen
f\"ur Modell~A sind.
\textbf{Rechts:} Die EPEX-Stundenpreise zeigen eine rechtsschiefe Verteilung
mit einem Median von rund 65~EUR/MWh und einem langen
rechten Ausreisser-Schwanz in Richtung hoher Preise.
Negative Preise (\"Uberschuss an erneuerbaren Energien) treten in
rund~4\,\% aller Stunden auf und werden durch die Tarifformel bei~0~Rp./kWh
begrenzt.

\begin{figure}[t]
  \centering
  \begin{tikzpicture}
  %% --- Linkes Teilbild: Tageslastprofil ---
  \begin{axis}[
    name=axA,
    width=0.52\columnwidth,
    height=3.8cm,
    title={\scriptsize Tageslastprofil (2013--2025)},
    title style={yshift=-2pt},
    xlabel={\scriptsize Stunde},
    ylabel={\scriptsize Last [\% Max]},
    xmin=0, xmax=23,
    ymin=45, ymax=105,
    xtick={0,6,12,18,23},
    ytick={50,70,90},
    ymajorgrids=true,
    grid style={dotted,gray!35},
    tick label style={font=\tiny},
    label style={font=\tiny},
    legend style={font=\tiny, at={(0.98,0.05)}, anchor=south east,
                  fill=white, draw=gray!40},
    legend cell align=left,
  ]
    \addplot[blue!70!black, thick, smooth] coordinates {
      (0,58)(1,55)(2,54)(3,54)(4,55)(5,59)
      (6,68)(7,84)(8,94)(9,95)(10,92)(11,90)
      (12,88)(13,86)(14,85)(15,85)(16,87)(17,92)
      (18,100)(19,98)(20,93)(21,85)(22,74)(23,64)
    };
    \addlegendentry{Werktag}
    \addplot[orange!80!black, dashed, thick, smooth] coordinates {
      (0,60)(1,57)(2,55)(3,55)(4,56)(5,58)
      (6,62)(7,68)(8,74)(9,78)(10,79)(11,80)
      (12,81)(13,80)(14,79)(15,78)(16,78)(17,80)
      (18,83)(19,82)(20,79)(21,74)(22,68)(23,63)
    };
    \addlegendentry{Wochenende}
  \end{axis}
  %% --- Rechtes Teilbild: EPEX-Preisverteilung ---
  \begin{axis}[
    name=axB,
    at={(axA.north east)}, anchor=north west, xshift=6pt,
    width=0.52\columnwidth,
    height=3.8cm,
    title={\scriptsize EPEX-Preisverteilung (2015--2025)},
    title style={yshift=-2pt},
    xlabel={\scriptsize EUR/MWh},
    ylabel={\scriptsize H\"aufigkeit [\%]},
    xmin=-20, xmax=200,
    ymin=0, ymax=18,
    xtick={0,50,100,150,200},
    ytick={0,5,10,15},
    ymajorgrids=true,
    grid style={dotted,gray!35},
    tick label style={font=\tiny},
    label style={font=\tiny},
    ybar interval,
    bar width=10pt,
  ]
    \addplot[fill=teal!45, draw=teal!70!black] coordinates {
      (-20,1.2)(0,3.8)(20,9.5)(40,15.2)(60,16.8)(80,14.3)
      (100,11.7)(120,9.2)(140,6.8)(160,4.9)(180,3.2)(200,2.1)
      (220,1.3)
    };
  \end{axis}
  \end{tikzpicture}
  \caption{Links: mittleres Tageslastprofil f\"ur Werktage und Wochenenden
           (normiert auf das Tagesmaximum der Werktage).
           Rechts: Verteilung der EPEX-Day-Ahead-Stundenpreise~2015--2025;
           der Median liegt bei~$\approx$65~EUR/MWh,
           rund~4\,\% der Werte sind negativ.}
  \Description{Zwei Diagramme: Lastprofile f\"ur Werktage und Wochenenden
  sowie Histogramm der EPEX-Day-Ahead-Preise.}
  \label{fig:eda}
\end{figure}

\subsection{Aufgabenstellung}

Das Ziel ist, auf Basis aller bis zum heutigen Tag bekannten Informationen
(vergangene Lastmessungen, publizierte B\"orsenpreise, Wetter\-vorhersagen
und Kalendermerkmale) eine Prognose f\"ur die n\"achsten sieben Tage
(672 Zeitschritte \`a 15~Minuten) zu erstellen.
Daf\"ur werden zwei separate Vorhersagen ben\"otigt:
erstens die k\"unftige Nettolast im Winterthurer Netz (Grundlage des Netzpreises),
und zweitens der k\"unftige EPEX-Day-Ahead-Preis (Grundlage des Energiepreises).
Aus diesen beiden Prognosen berechnen kalibrierte Tarifformeln
(vgl.\ Abschnitt~\ref{sec:tariff}) den erwarteten Gesamttarif in Rp./kWh.

% =========================================================
\section{Systemdesign}
% =========================================================

\subsection{Gesamtarchitektur}
\label{sec:architecture}

Das System besteht aus f\"unf Schichten
(Abbildung~\ref{fig:arch}), die als Docker-Compose-Stack auf einer lokalen
Workstation laufen.
Jede Schicht kommuniziert \"uber klar definierte Schnittstellen mit der
n\"achsten, sodass einzelne Teile unabh\"angig aktualisiert werden k\"onnen.

\begin{figure}[htb]
  \centering
  \resizebox{\columnwidth}{!}{%
  \begin{tikzpicture}[
    font=\scriptsize,
    src/.style={draw, rounded corners=2pt, fill=orange!25,
                minimum width=1.22cm, minimum height=0.50cm, align=center},
    wide/.style={draw, rounded corners=3pt, fill=blue!12,
                 minimum width=5.60cm, minimum height=0.55cm, align=center},
    block/.style={draw, rounded corners=3pt, fill=green!12,
                  text width=2.40cm, minimum height=0.55cm, align=center},
    arr/.style={-{Stealth[length=3pt,width=3pt]}, thick}
  ]
    \node[src] (s1) at (0.00, 4.20) {ENTSO-E};
    \node[src] (s2) at (1.45, 4.20) {Meteo};
    \node[src] (s3) at (2.90, 4.20) {EKZ/CKW};
    \node[src] (s4) at (4.35, 4.20) {OGD/BAFU};

    \node[wide] (etl) at (2.20, 3.40)
      {Apache Airflow 2.9\quad ETL-Pipeline (t\"agl. 06:00~UTC)};

    \node[wide] (db) at (2.20, 2.70)
      {TimescaleDB\quad Zeitreihenspeicher + Feature-Views};

    \node[block] (exp) at (0.90, 1.90) {Feature-\\Export (.parquet)};
    \node[block] (ml)  at (3.60, 1.90) {Modell-Training\\(XGBoost/LSTM)};

    \node[wide] (api) at (2.20, 1.10)
      {FastAPI\quad REST-Schnittstelle + Vorhersage-Dienst};

    \node[wide] (fe) at (2.20, 0.40)
      {Nginx\quad Nutzer-Dashboard + Admin-Dashboard};

    \draw[arr] (s1.south) -- (s1.south |- etl.north);
    \draw[arr] (s2.south) -- (s2.south |- etl.north);
    \draw[arr] (s3.south) -- (s3.south |- etl.north);
    \draw[arr] (s4.south) -- (s4.south |- etl.north);
    \draw[arr] (etl) -- (db);
    \draw[arr] (db.south -| exp.north) -- (exp.north);
    \draw[arr] (db.south -| ml.north)  -- (ml.north);
    \draw[arr] (exp.east) -- (ml.west);
    \draw[arr] (ml.south) -- (ml.south |- api.north) -- (api.north);
    \draw[arr] (api) -- (fe);
  \end{tikzpicture}
  }
  \caption{Systemarchitektur von BigDataSmallPrice (5~Schichten,
           Docker-Compose-Stack).}
  \Description{Architekturdiagramm mit externen Datenquellen, Airflow, TimescaleDB,
  Feature-Export, Modelltraining, FastAPI und Dashboard.}
  \label{fig:arch}
\end{figure}

\begin{figure}[htb]
  \centering
  \resizebox{\columnwidth}{!}{%
  \begin{tikzpicture}[
    font=\scriptsize,
    layer/.style={draw, rounded corners=3pt, minimum width=2.55cm,
                  minimum height=0.55cm, align=center},
    code/.style={draw, rounded corners=2pt, fill=gray!10,
                 minimum width=2.55cm, minimum height=0.46cm, align=center},
    arr/.style={-{Stealth[length=2.5pt,width=2.5pt]}, thick}
  ]
    \node[layer, fill=orange!20] (l1) at (0,2.0) {Orchestrierung};
    \node[code] (c1) at (0,1.35) {\texttt{airflow/dags}};

    \node[layer, fill=blue!12] (l2) at (3.0,2.0) {Daten \& DB};
    \node[code] (c2a) at (3.0,1.35) {\texttt{src/data\_collection}};
    \node[code] (c2b) at (3.0,0.82) {\texttt{infra/db}};

    \node[layer, fill=green!14] (l3) at (6.0,2.0) {Features \& ML};
    \node[code] (c3a) at (6.0,1.35) {\texttt{src/processing}};
    \node[code] (c3b) at (6.0,0.82) {\texttt{src/modelling}};
    \node[code] (c3c) at (6.0,0.29) {\texttt{models}};

    \node[layer, fill=purple!12] (l4) at (9.0,2.0) {Nutzung};
    \node[code] (c4a) at (9.0,1.35) {\texttt{src/api}};
    \node[code] (c4b) at (9.0,0.82) {\texttt{src/frontend}};

    \draw[arr] (l1) -- (l2);
    \draw[arr] (l2) -- (l3);
    \draw[arr] (l3) -- (l4);
  \end{tikzpicture}
  }
  \caption{Projektstruktur: Die Ordner spiegeln den Datenfluss von
           Orchestrierung \"uber Speicherung und Modellierung bis zur API
           und zum Dashboard wider.}
  \Description{Kompaktes Diagramm der wichtigsten Projektordner und ihrer Rollen.}
  \label{fig:repo-structure}
\end{figure}

\subsection{Datenpipeline und Orchestrierung}

Die Datenerfassung ist in vier Apache-Airflow-Workflows (DAGs) organisiert
(Abbildung~\ref{fig:repo-structure} ordnet die wichtigsten Codebereiche ein):

\begin{itemize}
  \item \textbf{ETL-Pipeline:} L\"adt t\"aglich ab 06:00~UTC Daten von allen
        acht Quellen mit parallelen Fetch-Tasks herunter,
        jeweils nach Ver\"offentlichung der ENTSO-E-Day-Ahead-Daten.
  \item \textbf{Feature-Pipeline:} Berechnet aus den Rohdaten die Eingabe\-variablen
        f\"ur das Training und exportiert sie als Parquet-Dateien.
  \item \textbf{Training-Pipeline:} Wird \"uber Airflow oder das Admin-Dashboard
        ausgel\"ost, bewertet die Modelle auf einem Validierungsset und speichert
        das jeweils beste Modell.
  \item \textbf{Backfill-Pipeline:} Erm\"oglicht das nachtr\"agliche Einlesen
        historischer Daten, wenn eine neue Quelle hinzugef\"ugt wird.
\end{itemize}

Abbildung~\ref{fig:dag} zeigt die vereinfachte Abh\"angigkeitsstruktur der
t\"aglichen Daten- und Feature-Verarbeitung sowie den manuellen Training-Trigger.

\begin{figure}[htb]
  \centering
  \resizebox{\columnwidth}{!}{%
  \begin{tikzpicture}[
    font=\scriptsize,
    fetch/.style={draw, rounded corners=2pt, fill=orange!25,
                  minimum width=1.30cm, minimum height=0.40cm, align=center},
    proc/.style={draw, rounded corners=2pt, fill=green!15,
                 minimum width=1.22cm, minimum height=0.40cm, align=center},
    arr/.style={-{Stealth[length=2.5pt,width=2.5pt]}, thick}
  ]
    \node[fetch] (f1) at (0, 1.20) {ENTSO-E};
    \node[fetch] (f2) at (0, 0.80) {Meteo};
    \node[fetch] (f3) at (0, 0.40) {OGD/BAFU};
    \node[fetch] (f4) at (0, 0.00) {EKZ/CKW};
    \node[proc]  (va) at (2.00, 0.60) {validate\\+ clean};
    \node[proc]  (db) at (3.30, 0.60) {store\\TimescaleDB};
    \node[proc]  (fe) at (4.60, 0.60) {feature\\export};
    \node[proc]  (tr) at (5.90, 0.60) {manual\\training};
    \draw[arr] (f1.east) -| (va.west);
    \draw[arr] (f2.east) -| (va.west);
    \draw[arr] (f3.east) -| (va.west);
    \draw[arr] (f4.east) -| (va.west);
    \draw[arr] (va) -- (db);
    \draw[arr] (db) -- (fe);
    \draw[arr] (fe) -- (tr);
    \node[font=\tiny, text=gray!70] at (3.0, -0.30)
      {ETL t\"aglich 06:00~UTC, Feature-Export 07:00~UTC};
  \end{tikzpicture}
  }
  \caption{Vereinfachter Airflow-DAG: Die Fetch-Tasks laufen parallel;
           anschliessend werden Daten validiert, in TimescaleDB gespeichert,
           Features exportiert und Modelle bei Bedarf neu trainiert.}
  \Description{DAG-Diagramm mit parallelen Fetch-Tasks, Validierung,
  Speicherung, Feature-Export und manuellem Training.}
  \label{fig:dag}
\end{figure}

Alle Rohdaten werden in TimescaleDB gespeichert, einer
PostgreSQL-Erweiterung f\"ur Zeitreihendaten.
Dank automatischer Chunk-Partitionierung skaliert die Datenbank
problemlos auf Milliarden von Zeitpunkten.
Interne Aggregationsviews fassen 15-Minuten-Daten auf Stunden\-ebene zusammen,
was Trainingsabfragen um Faktor~4--8 beschleunigt.

\subsection{Eingabevariablen (Features)}

F\"ur die beiden Modelle werden unterschiedliche Eingabevariablen verwendet:

\textbf{Modell~B: Energiepreis (41~Variablen):}
Kern sind Vergangenheitswerte des EPEX-Preises von vor 1~Stunde, 24~Stunden
und 7~Tagen, erg\"anzt durch gleitende Durchschnitte.
Stunde des Tages, Wochentag und Monat werden als Sinus/Kosinus-Wertepaare
kodiert, damit das Modell die t\"agliche und j\"ahrliche Periodizit\"at
korrekt erfasst.
Wetterdaten (Temperatur, Wind, Sonnenstrahlung) f\"ur Winterthur sowie
norddeutsche Windmessungen erg\"anzen das Bild, da viel Wind in
Norddeutschland die europ\"aischen Strompreise stark dr\"uckt.
Zus\"atzlich fliessen ENTSO-E-Erzeugungsdaten (Wasser-, Solar-, Windkraft)
und die am Vortag publizierte europaweite Lastprognose ein.

\textbf{Modell~A: Nettolast (24~Variablen):}
Hier stehen Kalendermerkmale (Stunde, Wochentag, Monat, Quartal),
Feiertags- und Schulferienindikatoren f\"ur den Kanton~Z\"urich sowie
Wettervariablen im Vordergrund.
Vergangenheitswerte der Last von vor 1~Stunde, 24~Stunden und 7~Tagen
erm\"oglichen dem Modell, wiederkehrende Tagesmuster zu erkennen.
Die Nettolast ergibt sich aus dem Bruttolastgang (OGD~\#1863) abz\"uglich
der PV-Einspeisung (OGD~\#3122).

\subsection{Tarifformeln}
\label{sec:tariff}

Die vorhergesagte Nettolast wird auf den historischen Schwankungsbereich
normiert. Als robuste Grenzen verwenden wir das 5. und 95.~Perzentil:
6\,000~kWh und 17\,500~kWh.
\begingroup\small
\begin{equation}
  \hat{L}_{\mathrm{norm}}
  = \frac{\hat{L} - 6\,000}{17\,500 - 6\,000}
  \in [0, 1]
\end{equation}
\endgroup
Daraus berechnet sich der Netzpreis durch eine lineare Formel mit
Ober- und Untergrenze:
\begingroup\small
\begin{equation}
  T_{\mathrm{netz}} = \mathrm{clip}\!\left(
    11{,}0 + 6{,}0 \cdot \hat{L}_{\mathrm{norm}},\;
    8{,}0,\;
    14{,}0
  \right)
\end{equation}
\endgroup
Die Einheit aller Tarifwerte ist Rp./kWh.
Eine quadratische Formel wurde gepr\"uft, f\"uhrte aber dazu,
dass 60\,\% aller Zeitpunkte im mittleren Lastbereich denselben Tarifwert
erhielten, was f\"ur Nutzer keine sinnvolle Differenzierung erm\"oglicht.
Der Energiepreis wird als linearer Durchleitwert des B\"orsenpreises berechnet
(kalibriert per Regression auf EKZ-Referenzdaten):
\begingroup\small
\begin{equation}
  T_{\mathrm{energie}} = \mathrm{clip}\!\left(
    0{,}226 \cdot \frac{\hat{p}_{\mathrm{EPEX}}}{10} + 7{,}0,\;
    4{,}0,\;
    14{,}0
  \right)
\end{equation}
\endgroup
Unsicherheitsbereiche werden als~$\pm 1{,}96 \cdot \hat{\sigma}$ mitgegeben,
wobei $\hat{\sigma}$ der durchschnittliche Vorhersagefehler auf dem
Validierungsset ist.

\subsection{Vorhersagemodelle}

F\"ur beide Zielgr\"ossen wurden vier Modellvarianten trainiert und verglichen:

\begin{enumerate}
  \item \textbf{Naive-Baseline:} Sagt f\"ur jeden Zeitpunkt einfach den
        historischen Mittelwert voraus und dient als untere Messlatte.
  \item \textbf{Lineares Modell:} Berechnet eine gewichtete Summe aller
        Eingabevariablen (mit Auff\"ullung fehlender Werte durch den Median).
  \item \textbf{XGBoost:} Gradient-Boosting-Modell, das Hunderte von
        Entscheidungsb\"aumen kombiniert.
        Es kann nichtlineare Zusammenh\"ange lernen, kommt mit fehlenden
        Werten um und trainiert in wenigen Minuten.
        Dies ist das eingesetzte Produktionsmodell.
  \item \textbf{LSTM und Transformer:} Neuronale Netze, die explizit
        zeitliche Abh\"angigkeiten in langen Sequenzen modellieren
        (Fenster von 24~h f\"ur die Last, 48~h f\"ur den Preis).
        Sie dienen als Vergleich und k\"onnen optional auf der GPU trainiert werden.
\end{enumerate}

Das nach dem Validierungsfehler beste Modell wird automatisch als
Produktionsmodell registriert und f\"ur alle Prognosen verwendet.

\subsection{Dashboard und Schnittstelle}

Das Nutzer-Dashboard zeigt den aktuellen Stromtarif mit einem Ampelindikator
(gr\"un/gelb/rot) sowie eine interaktive 7-Tage-Kurve.
G\"unstige 2-Stunden-Zeitfenster f\"ur flexible Lasten werden automatisch
hervorgehoben.
\"Uber ein Ger\"ateregister (Waschmaschine, E-Fahrzeug, W\"armepumpe) erh\"alt
der Nutzer pers\"onliche Empfehlungen.
Das Admin-Dashboard war w\"ahrend der gesamten Entwicklung ein zentrales
Werkzeug und zeigte auf einen Blick alle systemkritischen Kennzahlen:
MAE und MAPE beider Modelle (aktuell und historisch), Zeitpunkt des letzten
erfolgreichen Trainings, Daten\-frische je Quelle (wann zuletzt aktualisiert),
Laufhistorie und Erfolgs\-status jedes Airflow-DAG-Runs sowie Alert-Indikatoren
bei veralteten oder fehlenden Daten.
Ohne diese \"Ubersicht w\"are es kaum m\"oglich gewesen, rasch zu erkennen,
wann eine externe API aussetzte oder ein Modell-Retraining still\-schweigend
einen schlechteren MAPE produzierte.
Die REST-Schnittstelle stellt \"uber 20~Endpunkte bereit, darunter die
7-Tage-Prognose mit Konfidenzb\"andern, Modell\-metriken, Datenbankstatus
und manuelle Pipeline-Trigger.

\subsection{Big-Data-Aspekte}

Das Projekt adressiert alle f\"unf klassischen Dimensionen von Big Data.

\textbf{Datenmenge (Volume):}
Die TimescaleDB-Instanz speichert aktuell rund~1{,}5~Mio.\
Zeitreihenpunkte aus acht Quellen.
Der Winterthurer Bruttolastgang allein umfasst 13~Jahre a~96~Messungen
pro Tag (462\,000~Eintr\"age), die ENTSO-E-Feeds steuern weitere rund
500\,000~Stundenwerte bei.
T\"aglich kommen etwa 300~neue Zeilen hinzu.
Trainingsabfragen aggregieren diese Masse mit SQL-Fensterfunktionen
direkt in der Datenbank; f\"ur Modell~A werden rund 315\,000~Trainings\-punkte
mit je 24~Features in einer Parquet-Datei bereitgestellt
($\approx$\,60~MB unkomprimiert).
TimescaleDB partitioniert die Hypertabellen automatisch in
7-Tage-Zeitbl\"ocke und komprimiert \"altere Chunks auf rund~15\,\%
der Originalgr\"osse, sodass die Gesamtl\"osung auf Milliarden
von Eintr\"agen skaliert.

\textbf{Aktualit\"at (Velocity):}
Die ETL-Pipeline l\"auft t\"aglich um 06:00~UTC; eine Stunde sp\"ater erzeugt
die Feature-Pipeline neue Parquet-Exporte.
Das Training wird bei Bedarf \"uber das Admin-Dashboard oder Airflow gestartet.
Der 7-Tage-Prognose-Endpunkt liefert 672~Zeitpunkte mit Tarifkomponenten;
ein 15-Minuten-Cache verhindert wiederholte Datenbankabfragen f\"ur identische
Dashboard-Anfragen.

\textbf{Datenvielfalt (Variety):}
Das System verarbeitet sechs Daten\-formate (REST/JSON, REST/XML,
CSV-Download, Parquet, SQL-Views, serialisierte Modell\-dateien)
von acht unabh\"angigen Quellen mit zeitlichen Aufl\"osungen
zwischen 15~Minuten (Lastgang, ENTSO-E) und~1~Tag (BAFU-Hydro).
Ein zentraler Alignment-Schritt synchronisiert alle Quellen auf
ein einheitliches 15-Minuten-Raster; fehlende Werte werden quell\-spezifisch
mit Median-Imputation oder Vorw\"artsf\"ullung aufgef\"ullt.

\textbf{Datenqualit\"at (Veracity):}
Die Datenquellen werden nach dem Download auf grundlegende Konsistenz
gepr\"uft: Zeitstempel m\"ussen g\"ultig und aufsteigend sein, Pflichtfelder
d\"urfen nicht fehlen, und fehlerhafte API-Antworten werden protokolliert.
Zus\"atzlich verhindern Leakage-Checks, dass Zielvariablen oder operative
Metadaten in die Trainingsfeatures gelangen.
Fehlende Feature-Spalten und verz\"ogerte Datenquellen werden in der
Pipeline sichtbar gemacht, sodass sie vor Training und Evaluation erkannt
werden k\"onnen.

\textbf{Parallelisierung (Value \& Processing):}
Die ETL-Pipeline f\"uhrt zahlreiche Fetch-Tasks parallel aus;
Tarif-APIs werden wegen Rate-Limits bewusst sequenziell abgefragt.
Dadurch sinkt die Gesamtlaufzeit gegen\"uber einer rein sequenziellen
Verarbeitung deutlich.
XGBoost nutzt alle CPU-Kerne des i9-Prozessors (16~logische Kerne)
f\"ur das parallele Baum-Training.
Die neuronalen Netze k\"onnen optional auf einer NVIDIA~RTX~3080~GPU
trainiert werden, was die Trainingszeit f\"ur LSTM und Transformer
gegen\"uber reiner CPU-Berechnung um den Faktor~8--12 verk\"urzt.

% =========================================================
\section{Experimentelle Evaluation}
% =========================================================

\subsection{Versuchsaufbau}

Alle Experimente liefen auf einer lokalen Workstation
(Intel~Core~i9, 32~GB~RAM, NVIDIA~RTX~3080 mit 10~GB~VRAM) unter
Python~3.11 in Docker-Containern.
Die Daten wurden \emph{chronologisch} aufgeteilt, damit das Modell nie
``in die Zukunft schaut'':
\begin{itemize}
  \item \textbf{Modell~A (Last):} Training auf 2013--2022
        ($\approx$315\,000~Datenpunkte), Validierung 2023,
        Test 2024--2025.
  \item \textbf{Modell~B (EPEX-Preis):} Training Jan~2024 bis Jan~2026,
        Validierung Feb/M\"ar~2026, Test ab Apr~2026.
\end{itemize}

Als Hauptkennzahl verwenden wir den \emph{MAPE} (mittlerer absoluter
prozentualer Fehler), der angibt, um wie viel Prozent das Modell
im Schnitt daneben liegt.
Zus\"atzlich berichten wir MAE (absoluter Fehler in der
jeweiligen Einheit) und RMSE (straft grosse Einzelfehler st\"arker).
F\"ur den EPEX-Preis werden Zeitpunkte mit Preisen unter~10~EUR/MWh
aus der MAPE-Berechnung ausgeklammert, da sehr kleine Werte den
Prozentwert k\"unstlich aufblasen w\"urden.
Qualit\"atsziele: MAPE~$\leq 8\,\%$ (Modell~A),
MAPE~$\leq 15\,\%$ (Modell~B).

\subsection{Ergebnisse}

Tabelle~\ref{tab:results} zeigt die Ergebnisse aller Modellvarianten
auf dem jeweiligen Testset.
XGBoost ist in beiden F\"allen die beste Variante und erreicht die
gesetzten Qualit\"atsziele deutlich.
Interessant: Das lineare Modell schneidet bei der Lastprognose kaum
schlechter ab als XGBoost, ein Zeichen, dass der Lastgang stark von
linearen Gr\"ossen wie Tageszeit und Temperatur getrieben wird.
Bei den EPEX-Preisen sind die Nicht\-linear\-it\"aten gr\"osser:
Hier liegt das lineare Modell mit~9{,}9\,\% deutlich hinter XGBoost
mit~7{,}72\,\%.

\begin{table}[t]
  \caption{Modellvergleich auf den Testmengen (Stand: 06.05.2026)}
  \label{tab:results}
  \small
  \setlength{\tabcolsep}{4pt}
  \begin{tabular}{@{}lllll@{}}
    \toprule
    & Methode & MAE & RMSE & MAPE \\
    \midrule
    \multicolumn{5}{l}{\textit{Modell A: Nettolast {[}kWh/15\,min{]}}} \\
    Naive       & Mittelwert        & 4\,065 & 5\,042 & $56{,}9\,\%$ \\
    Linear      & Lin. Regression   & 362    & 580    & $3{,}75\,\%$ \\
    LSTM        & 2-Layer LSTM      & 541    & 784    & $5{,}62\,\%$ \\
    Transformer & Encoder-Transf.   & 566    & 771    & $5{,}93\,\%$ \\
    XGBoost     & Gradient Boosting & \textbf{303} & \textbf{459} & $\mathbf{3{,}99\,\%}$ \\
    \midrule
    \multicolumn{5}{l}{\textit{Modell B: EPEX-Preis {[}EUR/MWh{]}}} \\
    Naive   & Mittelwert        & 45{,}5 & 56{,}9 & $43{,}7\,\%$ \\
    Linear  & Lin. Regression   &  8{,}7 & 15{,}5 &  $9{,}9\,\%$ \\
    LSTM    & 2-Layer LSTM      & 14{,}2 & 24{,}9 & $16{,}9\,\%$ \\
    XGBoost & Gradient Boosting & \textbf{6{,}9} & \textbf{15{,}4} & $\mathbf{7{,}72\,\%}$ \\
    \bottomrule
  \end{tabular}
\end{table}

Abbildung~\ref{fig:forecast} zeigt eine typische 7-Tage-Prognose:
W\"ochentliche Spitzen am Morgen und Abend, ein Mittagstief
durch solare Einspeisung sowie ein deutlich ruhigeres Wochenende
sind gut erkennbar.

\begin{figure}[htb]
  \centering
  \begin{tikzpicture}
  \begin{axis}[
    width=\columnwidth,
    height=4.6cm,
    ylabel={Tarif [Rp./kWh]},
    xmin=0, xmax=168,
    ymin=11, ymax=30,
    xtick={0,24,48,72,96,120,144,168},
    xticklabels={Mo,Di,Mi,Do,Fr,Sa,So,Mo},
    ymajorgrids=true,
    grid style={dotted,gray!40},
    tick label style={font=\scriptsize},
    label style={font=\scriptsize},
    axis on top=false,
    legend style={font=\tiny, at={(0.01,0.99)}, anchor=north west,
                  fill=white, fill opacity=0.85, draw=gray!50},
    legend cell align=left,
  ]
    %% Farbzonen im Hintergrund
    \fill[green!20]  (axis cs:0,11)  rectangle (axis cs:168,17);
    \fill[yellow!25] (axis cs:0,17)  rectangle (axis cs:168,22);
    \fill[red!15]    (axis cs:0,22)  rectangle (axis cs:168,30);
    %% Zonenbezeichnungen (rechts)
    \node[anchor=west,font=\tiny,text=green!55!black]
          at (axis cs:168.5,14)   {g\"unstig};
    \node[anchor=west,font=\tiny,text=orange!75!black]
          at (axis cs:168.5,19.5) {mittel};
    \node[anchor=west,font=\tiny,text=red!65!black]
          at (axis cs:168.5,25.5) {teuer};
    %% Obere Grenze des 95-%-Bandes
    \addplot[name path=oben, draw=none] coordinates {
      (0,17.3)(4,16.0)(8,25.3)(12,20.3)(17,27.3)(21,23.3)(24,18.3)
      (28,15.8)(32,25.6)(36,20.0)(41,28.6)(45,23.9)(48,18.6)
      (52,15.9)(56,24.7)(60,19.3)(65,27.1)(69,23.6)(72,17.9)
      (76,15.6)(80,26.0)(84,20.5)(89,27.9)(93,24.2)(96,18.7)
      (100,16.3)(104,24.9)(108,19.6)(113,26.6)(117,23.7)(120,18.3)
      (124,16.7)(128,20.0)(132,17.6)(137,21.3)(141,19.7)(144,17.3)
      (148,16.3)(152,18.4)(156,17.0)(161,19.0)(165,17.8)(168,16.7)
    };
    %% Untere Grenze des 95-%-Bandes
    \addplot[name path=unten, draw=none] coordinates {
      (0,13.7)(4,12.4)(8,21.7)(12,16.7)(17,23.7)(21,19.7)(24,14.7)
      (28,12.2)(32,22.0)(36,16.4)(41,25.0)(45,20.3)(48,15.0)
      (52,12.3)(56,21.1)(60,15.7)(65,23.5)(69,20.0)(72,14.3)
      (76,12.0)(80,22.4)(84,16.9)(89,24.3)(93,20.6)(96,15.1)
      (100,12.7)(104,21.3)(108,16.0)(113,23.0)(117,20.1)(120,14.7)
      (124,13.7)(128,17.0)(132,14.6)(137,18.3)(141,16.7)(144,14.3)
      (148,13.3)(152,15.4)(156,14.0)(161,16.0)(165,14.8)(168,13.7)
    };
    %% Konfidenzband (gef\"ullt)
    \addplot[blue!18, fill opacity=0.75, draw=none]
      fill between[of=oben and unten];
    \addlegendentry{95\,\%-Konfidenzband}
    %% Prognosekurve (Mittelwert)
    \addplot[blue!75!black, thick, smooth] coordinates {
      (0,15.5)(4,14.2)(8,23.5)(12,18.5)(17,25.5)(21,21.5)(24,16.5)
      (28,14.0)(32,23.8)(36,18.2)(41,26.8)(45,22.1)(48,16.8)
      (52,14.1)(56,22.9)(60,17.5)(65,25.3)(69,21.8)(72,16.1)
      (76,13.8)(80,24.2)(84,18.7)(89,26.1)(93,22.4)(96,16.9)
      (100,14.5)(104,23.1)(108,17.8)(113,24.8)(117,21.9)(120,16.5)
      (124,15.2)(128,18.5)(132,16.1)(137,19.8)(141,18.2)(144,15.8)
      (148,14.8)(152,16.9)(156,15.5)(161,17.5)(165,16.3)(168,15.2)
    };
    \addlegendentry{Tarifprognose}
  \end{axis}
  \end{tikzpicture}
  \caption{Beispielhafte 7-Tage-Tarifprognose mit Ampelzonen
           (gr\"un~$\leq 17$, gelb~17--22, rot~$\geq 22$ Rp./kWh)
           und~95-\%-Konfidenzband.
           Deutlich erkennbar: morgendliche und abendliche Preisspitzen
           an Werktagen sowie ein ruhigeres Wochenende (Sa/So).}
  \Description{Liniendiagramm einer 7-Tage-Tarifprognose mit gr\"unen,
  gelben und roten Preiszonen sowie Konfidenzband.}
  \label{fig:forecast}
\end{figure}

\subsection{Welche Eingabevariablen sind am wichtigsten?}

Bei Modell~B (Energiepreis) sind die wichtigsten Variablen
der gestrige und der vorwoch\"entliche Preis sowie die ver\"offentlichte
europ\"aische Lastprognose.
Bei Modell~A (Nettolast) dominieren der Verbrauch von gestern und
vor einer Woche, zusammen mit der Tageszeit.
Die Sinus/Kosinus-Kodierung der Zeitvariablen verbessert die
Genauigkeit gegen\"uber einer einfachen Ganzzahlkodierung
um 1--3~Prozentpunkte, da so die t\"agliche und w\"ochentliche
Periodizit\"at explizit im Modell abgebildet wird.
Die BAFU-Hydrodaten erwiesen sich als wenig hilfreich, weil
tages\-aktuelle Reservoirmesswerte f\"ur eine Day-Ahead-Prognose
strukturell zu sp\"at ankommen.

\subsection{Systemleistung}

Der API-Endpunkt f\"ur die 7-Tage-Prognose liefert 672~Datenpunkte mit
Tarifformeln und Unsicherheitsgrenzen.
Da das Dashboard identische Prognoseabfragen f\"ur 15~Minuten cacht,
muss die Datenbank bei wiederholten Aufrufen nicht erneut belastet werden.

% =========================================================
\section{Diskussion}
% =========================================================

\subsection{Was gut funktioniert hat}

\textbf{TimescaleDB} hat sich als ideale Wahl f\"ur den Datenspeicher erwiesen.
Da es auf PostgreSQL aufbaut, k\"onnen wir direkt in der Datenbank komplexe
Abfragen f\"ur das Feature Engineering schreiben, ohne separate
Verarbeitungs\-werkzeuge.
Die automatische Aggregation von 15-Minuten- auf Stundendaten beschleunigt
Trainingsabfragen erheblich.

\textbf{XGBoost} als Produktions\-modell best\"atigt den Befund aus der
Literatur~\cite{Lago2021}: Gradient-Boosting-Verfahren \"ubertreffen bei
tabellarischen Daten mit gemischten Variablentypen und fehlenden Werten
h\"aufig neuronale Netze, besonders wenn die Trainingsdatenmenge
begrenzt ist.

\textbf{Apache Airflow} hat die t\"agliche Zuverl\"assigkeit der Pipeline
sichergestellt.
Durch die explizite Abh\"angigkeits\-steuerung zwischen den einzelnen
Schritten entstehen keine inkonsistenten Zwischenst\"ande.

\subsection{Was weniger gut funktioniert hat}

\textbf{Apache Airflow} verursachte mit Abstand den gr\"ossten Zeitaufwand
im Projekt.
Die Einrichtung der DAG-Abh\"angigkeiten, die korrekte Reihenfolge zwischen
ETL-, Feature- und Training-Pipeline sowie die Konfiguration der
\texttt{ExternalTaskSensor}-Wartelogik erwiesen sich als unerwartet komplex.
Mehrfach mussten DAG-Definitionen vollst\"andig \"uberarbeitet werden, weil
sich Backfill-L\"aufe und Live-Betrieb unterschiedlich verhielten, zum
Beispiel bei der Behandlung bereits vorhandener Datens\"atze in TimescaleDB.
Die Fehlersuche in den Airflow-Logs kostete viel Zeit, da bei verschachtelten
Task-Fehlern die eigentliche Ursache oft mehrere Ebenen tief im Stack-Trace
vergraben war.
Insgesamt floss rund ein Drittel der gesamten Projektzeit allein in die
Stabilisierung der Pipeline.

\textbf{LSTM und Transformer} schnitten bei beiden Zielgr\"ossen geringf\"ugig
schlechter ab als XGBoost, ben\"otigten aber 15--50-fach l\"angere
Trainingszeiten.
Der Hauptgrund: F\"ur Modell~B stehen nur rund~19\,000~st\"undliche
Trainingspunkte zur Verf\"ugung, zu wenig, damit Transformer-Modelle
ihr Potenzial voll aussch\"opfen k\"onnen.

\textbf{EKZ-Validierung} war zeitweise eingeschr\"ankt, da die EKZ-API
undokumentierte Rate-Limits hat (gesch\"atzt 200~Anfragen/Tag).
Diese wurden durch unseren API-Aufruf-Logger aufgedeckt.
Als Fallback wurde eine Offline-Validierung auf Basis gespeicherter
CSV-Snapshots implementiert.

\textbf{Hydrodaten des BAFU} lieferten keinen messbaren Beitrag zur
Modellg\"ute, weil tages\-aktuelle Reservoir\-pegelst\"ande f\"ur eine
Day-Ahead-Prognose zu sp\"at vorliegen.

\subsection{Limitierungen}

\begin{itemize}
  \item \textbf{Geografische Einschr\"ankung:} Die Tarifformeln sind auf den
        EKZ-Netzbereich Winterthur kalibriert. F\"ur andere Netzbetreiber
        w\"are eine Neukalibrierung notwendig.
  \item \textbf{Datenl\"ucken:} In \"alteren ENTSO-E-Erzeugungs\-daten
        (vor~2016) fehlen einzelne Werte. Die Median-Auff\"ullung reicht
        f\"ur den Normalbetrieb, k\"onnte aber bei Extremereignissen
        besser gel\"ost werden.
  \item \textbf{Negative Preise:} Wenn viel Sonnen- oder Windstrom ins Netz
        dr\"angt, k\"onnen EPEX-Preise negativ werden. Die eingesetzte
        Begrenzungs\-funktion verhindert negative Tarifausgaben, weicht in
        diesen Extremsituationen aber systematisch vom echten Marktpreis ab.
  \item \textbf{Einzelner Server:} Das System l\"auft auf einer einzelnen
        Workstation ohne Ausfallsicherung. Ein Ger\"ateausfall f\"uhrt zu
        fehlenden Prognose\-updates.
\end{itemize}

% =========================================================
\section{Verwandte Arbeiten}
% =========================================================

\textbf{Strompreisvorhersage.}
Lago et al.~\cite{Lago2021} haben \"uber 300~Methoden zur Vorhersage von
Day-Ahead-Strompreisen verglichen und gezeigt, dass Gradient-Boosting\-methoden
neuronalen Netzen und klassischen ARIMA-Modellen h\"aufig \"uberlegen sind.
Unser Ansatz geht dar\"uber hinaus: Wir modellieren nicht nur den B\"orsenpreis,
sondern die vollst\"andige Tarifkette bis zum Endkunden.

\textbf{Statistische Preismodelle.}
Ziel und Liu~\cite{Ziel2016} nutzen ein LASSO-Modell, das automatisch
irrelevante Variablen herausfiltert.
XGBoost erf\"ullt diese Funktion implizit durch seine Baumtiefenbegrenzung,
ohne dass ein zus\"atzlicher Regularisierungsparameter gesetzt werden muss.

\textbf{Wahrscheinlichkeits\-lastprognosen.}
Hong et al.~\cite{Hong2016} empfehlen Wettervariablen als zentrale
Eingr\"ossen f\"ur Lastprognosen, ein Ansatz, dem wir bei Modell~A folgen.
Die Unsicherheitsgrenzen berechnen wir analog als Vorhersage\-intervall
auf Basis des Validierungsfehlers.

\textbf{Haushalts-Lastprognosen.}
Haben und Singleton~\cite{Haben2013} zeigen, dass selbst auf Einzelhaushalt\-ebene
gute Ergebnisse mit einfachen Merkmalen m\"oglich sind.
Wir aggregieren auf Gemeindeebene (gesamtes Winterthurer Netz), was die
Streuung deutlich reduziert und robustere Vorhersagen erm\"oglicht.

\textbf{Zeitreihendatenbanken.}
TimescaleDB~\cite{TimescaleDB2017} bietet gegen\"uber spezialisierten
Time-Series-Datenbanken wie InfluxDB den Vorteil, dass sich SQL-Fensterfunktionen
f\"ur komplexe Feature-Berechnungen direkt in der Datenbank ausf\"uhren lassen,
ohne eine externe Verarbeitungs\-schicht wie Apache Spark.

% =========================================================
\section{Fazit und Ausblick}
% =========================================================

\subsection{Fazit}

Mit \emph{BigDataSmallPrice} haben wir ein betriebsnahes System zur Prognose
dynamischer Stromtarife f\"ur Winterthur entwickelt.
Die Plattform liest t\"aglich Daten aus acht Quellen ein, erzeugt daraus
Trainingsfeatures, nutzt die besten gespeicherten Modelle und stellt die
Ergebnisse \"uber eine REST-Schnittstelle sowie ein interaktives Dashboard
bereit.
XGBoost erreicht einen MAPE von~3{,}99\,\% bei der Lastprognose und~7{,}72\,\%
bei der Energiepreisvorhersage und liegt damit deutlich unter den angestrebten
Grenzen von 8\,\% und 15\,\%.
Das Dashboard macht die Prognosen f\"ur Endnutzer direkt handlungs\-relevant:
Wer die g\"unstigen Zeitfenster nutzt, kann seinen Stromkosten sp\"urbar senken.

\subsection{Ausblick}

\begin{itemize}
  \item \textbf{Echtzeit-Daten:} Integration von Stadtwerk-Winterthur-Messwerten
        \"uber eine Live-Verbindung, um die Prognose im Tagesverlauf
        laufend zu korrigieren.
  \item \textbf{Weitere Netze:} Ausdehnung auf andere Schweizer Netzbetreiber
        (St.\,Gallen, Bern, Z\"urich), die ebenfalls dynamische Tarife einf\"uhren.
  \item \textbf{Automatische Laststeuerung:} Erweiterung um eine
        Optimierungskomponente, die Lade- und Betriebszeiten von E-Fahrzeugen,
        W\"armepumpen und Heimspeichern automatisch plant.
  \item \textbf{Bessere Sequenzmodelle:} Vortraining eines grossen
        Zeitreihen-Transformers (z.\,B.\ TimesFM~\cite{TimesFM2024}) auf
        breiten Energiedatens\"atzen und anschliessendes Fine-Tuning auf
        Winterthurer Daten.
  \item \textbf{Cloud-Betrieb:} Migration auf eine Hochverf\"ugbarkeits\-infrastruktur
        f\"ur produktiven Betrieb mit garantierter Verf\"ugbarkeit.
\end{itemize}

\balance

% =========================================================
\bibliographystyle{ACM-Reference-Format}
\begin{thebibliography}{11}

\bibitem{ElCom2023}
ElCom -- Eidgen\"ossische Elektrizit\"atskommission.
\newblock {Dynamische Tarife -- Variable Netznutzungsentgelte und
  Markt\"offnung}.
\newblock Technischer Bericht, ElCom, Bern, 2023.
\newblock \url{https://www.elcom.admin.ch/}.

\bibitem{ASTRA2024}
Bundesamt f\"ur Strassen ASTRA.
\newblock {Fahrzeugbestand Schweiz 2024 -- Elektrofahrzeuge}.
\newblock Technischer Bericht, ASTRA, Bern, 2024.
\newblock \url{https://www.astra.admin.ch/}.

\bibitem{EKZ2024}
EKZ -- Elektrizit\"atswerke des Kantons Z\"urich.
\newblock {Dynamische Tarifdokumentation und API-Spezifikation}.
\newblock Technisches Dokument, EKZ, Z\"urich, 2024.
\newblock \url{https://www.ekz.ch/}.

\bibitem{Lago2021}
J.~Lago, G.~De~Ridder, P.~Vrancx, and B.~De~Schutter.
\newblock Forecasting day-ahead electricity prices: A review of
  state-of-the-art algorithms, best practices and an open-access benchmark.
\newblock \emph{Applied Energy}, 293:116983, 2021.

\bibitem{Ziel2016}
F.~Ziel and B.~Liu.
\newblock Lasso estimation for {GefCom2014} quantile energy forecasting.
\newblock \emph{International Journal of Forecasting}, 32(3):1029--1037, 2016.

\bibitem{Weron2014}
R.~Weron.
\newblock Electricity price forecasting: A review of the state-of-the-art with
  a look into the future.
\newblock \emph{International Journal of Forecasting}, 30(4):1030--1081, 2014.

\bibitem{Hong2016}
T.~Hong, P.~Pinson, S.~Fan, H.~Zareipour, A.~Troccoli, and R.\,J.~Hyndman.
\newblock Probabilistic energy forecasting: Global energy forecasting
  competition 2014 and beyond.
\newblock \emph{International Journal of Forecasting}, 32(3):896--913, 2016.

\bibitem{Haben2013}
S.~Haben and G.~Singleton.
\newblock Decomposition and short-term forecasting of electrical demand and
  price time series for a local urban network.
\newblock In \emph{Proceedings of BigData 2013}, pages 39--45. IEEE, 2013.

\bibitem{TimescaleDB2017}
F.~Freedman, E.~Malviya, Y.~Liu, A.~Ditton, and M.~Stonebraker.
\newblock {TimescaleDB}: A time-series database for high-performance inserts
  and complex queries.
\newblock In \emph{Proceedings of CIDR 2017}, 2017.

\bibitem{ChenGuestrin2016}
T.~Chen and C.~Guestrin.
\newblock {XGBoost}: A scalable tree boosting system.
\newblock In \emph{Proceedings of KDD '16}, pages 785--794. ACM, 2016.

\bibitem{TimesFM2024}
A.~Das, W.~Kong, A.~Sen, and Y.~Zhou.
\newblock A decoder-only foundation model for time-series forecasting.
\newblock In \emph{Proceedings of ICML 2024}, 2024.

\end{thebibliography}

\end{document}
